\section{Preamble}
\label{sec:flow:preamble}

\subsection{Run-time and compile-time evaluation}
\label{sec:flow:cte}

Control flow is normally resolved at run time: the generated machine code carries
the branch instructions, and which branch is taken depends on values the program
only has while running.

The keyword \token{cte} -- for \textit{Compile Time Execution} -- moves that
decision to compile time, for the two constructs that accept it: \token{cte if}
(Section~\ref{sec:flow:cte_if}) and \token{cte for}
(Section~\ref{sec:flow:cte_for}). What it demands in exchange is that the
compiler already know the values the decision rests on. Control flow evaluated
this way is the same mechanism as conditional compilation, the subject of
Chapter~\ref{chap:conditional_compilation}.

The converse also holds. A run-time test whose outcome the compiler knows is
rejected, because the branch it guards is either always taken or dead code. The
rule applies to the condition of \token{if} and \token{while}, and to the
pattern of \token{if let} and \token{while let}; each section lists the
diagnostic it raises.

\subsection{Unreachable code}
\label{sec:flow:unreachable}

\token{return}, \token{throw}, \token{panic}, \token{break} and
\token{continue} each transfer control out of the block that contains them. An
instruction that follows one of them in the same block can never run, and is an
error, E4260 ``unreachable statement''.

\begin{lstlisting}[style=coloredverbatimError, escapechar=@]
let mut i = 0;
loop {
  i += 1;
  break;
  @\hb{println(i);}@ // error, unreachable statement
}
\end{lstlisting}

A branch that ends with one of these five does not produce a value, and is
called a \textit{breaking branch}. The value rules of the constructs below
ignore breaking branches: a conditional or a match whose other branches all
produce a value still has one.
